\subsubsection{Supersonic Outflow Boundary Condition}
\label{SEC:BC:Outflow:Supersonic}

This is the reverse of the supersonic inflow boundary condition.
%
Since the outgoing flow is supersonic, no information from the boundary
condition can propagate upstream into the computational domain.
%
Therefore, the right state is simply fixed to the interior\slash left state,
\thatis
%
\begin{equation*}
  \cvVecR = \cvVecL
\end{equation*}

%%%%%%%%%%%%%%%%%%%%%%%%%%%%%%%%%%%%%%%%%%%%%%%%%%%%%%%%%%%%%%%%%%%%%%%%%%%%%%%%

If the outflow is adjacent to a no-slip wall, this boundary condition could
result in stability problems within the subsonic region of the boundary layer.
%
If stability issues occur, it is recommended to switch to either a
characteristic (\ref{SEC:BC:Farfield:Characteristic}) or
general outflow (\ref{SEC:BC:Outflow:General}) boundary condition, since both of
these boundary conditions test the outgoing Mach number to determine whether to
use a subsonic or supersonic boundary condition.

%###############################################################################

